\documentclass[11pt]{article}

\usepackage[margin=1in]{geometry}
\usepackage{amsmath,amssymb,amsthm}
\usepackage{enumitem}
\usepackage[hidelinks]{hyperref}

\title{From lateral parabolic Schiffer to classical Schiffer:\\
a spectral reduction}
\author{}
\date{\today}

\newcommand{\R}{\mathbb{R}}
\newcommand{\Om}{\Omega}
\newcommand{\eps}{\varepsilon}
\newcommand{\1}{\mathbf 1}
\newcommand{\nor}{\partial_\nu}
\newcommand{\tr}{\mathrm{tr}}

\newtheorem{lem}{Lemma}
\newtheorem{prop}[lem]{Proposition}
\newtheorem{cor}[lem]{Corollary}
\newtheorem{thm}[lem]{Theorem}
\newtheorem{rem}[lem]{Remark}

\begin{document}
\maketitle

\section{Statement}

The previous note (\texttt{direction2-heat-extension-computation.tex})
showed that the lateral parabolic Schiffer condition on a weighted
pair $(\Om,V)$ is strictly stronger than stationary weighted Schiffer,
and identified the second-order obstruction
\(
    \eta_2|_{\partial\Om} = 2\tr(\nabla^2V\,\nabla^2\phi)
    +(\nabla V)^\top\nabla^2\phi(\nabla V).
\)
The natural next step was to compute the higher residues $\eta_k$ and
test what the joint vanishing imposes. The spectral approach below
sums the entire tower at once and produces the following clean
reduction.

\begin{thm}\label{thm:main}
Let $\Om\subset\R^n$ be a smooth bounded domain, $V\in C^\infty(\overline\Om)$,
and $(\Om,V,\phi,\lambda)$ a smooth weighted Schiffer pair:
\[
    L_{\Om,V}\phi=\lambda\phi\ \text{in }\Om,\quad
    \nor\phi=0\ \text{on }\partial\Om,\quad
    \phi|_{\partial\Om}=c,\quad \phi\not\equiv 0.
\]
Let $h=h_{\Om,V}$ be the heat extension on $\Om\times[0,1]$ defined by
(2.11) of arXiv:2412.06344 (unweighted Laplacian, Neumann lateral
data, initial data $\phi$). If
\[
    h(x,t)\equiv c\qquad\text{for all }(x,t)\in\partial\Om\times[0,1],
\]
then $\lambda$ is a Neumann eigenvalue of the unweighted Laplacian
$-\Delta$ on $\Om$, and the unweighted Neumann spectral projection
$u:=P_\lambda\phi$ is a nontrivial classical Schiffer overdetermined
eigenfunction:
\[
    -\Delta u=\lambda u\ \text{in }\Om,\qquad
    \nor u=0\ \text{on }\partial\Om,\qquad
    u|_{\partial\Om}=c.
\]
\end{thm}

The conclusion is striking: lateral parabolic Schiffer for the
\emph{weighted} pair $(\Om,V)$ implies classical Schiffer for the
\emph{unweighted} Laplacian on the \emph{base} domain $\Om$ alone,
with the potential $V$ playing no role in the limit object beyond
having selected an eigenvalue $\lambda$ that lies in both spectra.

\section{Proof}

Throughout, $\Om$ is bounded with smooth boundary, $V$ smooth on
$\overline\Om$, and we use the unweighted Neumann Laplacian
$\Delta_N$ on $\Om$ with $L^2(dx)$-orthonormal eigenbasis
$\{\psi_j\}_{j\ge 0}$ and eigenvalues
$0=\mu_0<\mu_1\le\mu_2\le\cdots\to\infty$, so $-\Delta\psi_j=\mu_j\psi_j$
and $\nor\psi_j=0$ on $\partial\Om$.

Since $\Om$ is bounded and $V\in C^\infty(\overline\Om)$, the measures
$dx$ and $e^{-V}dx$ are equivalent: $L^2(dx)=L^2(e^{-V}dx)$ as sets
with equivalent norms. The eigenfunction $\phi$ admits a unique
$L^2(dx)$ expansion
\[
    \phi=\sum_{j\ge 0}\alpha_j\psi_j,\qquad
    \alpha_j=\int_\Om\phi\,\psi_j\,dx.
\]

\paragraph{Step 1: explicit heat extension.}
The bounded solution of
\[
    \partial_t h=\tfrac{1}{8}(\Delta_x h+\lambda h),\quad
    \nor h=0,\quad h(x,0)=\phi(x)
\]
on $\Om\times(0,1)$ is, by separation of variables in the Neumann
basis,
\begin{equation}\label{eq:hseries}
    h(x,t)=\sum_{j\ge 0}\alpha_j\,e^{t(\lambda-\mu_j)/8}\,\psi_j(x).
\end{equation}
The series converges in $C^\infty(\overline\Om\times[0,1])$ because
$\phi\in C^\infty(\overline\Om)$ (so $\alpha_j$ decay faster than any
power), at most finitely many $\mu_j$ are smaller than $\lambda$, and
the rest produce exponentially decaying factors.

\paragraph{Step 2: collapse on the lateral boundary.}
Restrict \eqref{eq:hseries} to $x\in\partial\Om$ and group terms by
distinct Neumann eigenvalue. Let $\Sigma:=\{\mu_j:j\ge 0\}$ be the set
of distinct Neumann eigenvalues of $-\Delta$ on $\Om$, and for each
$\mu\in\Sigma$ let
\[
    f_\mu(x):=\sum_{j:\mu_j=\mu}\alpha_j\,\psi_j(x)|_{\partial\Om}
        =(P_\mu\phi)(x)\big|_{\partial\Om},
\]
where $P_\mu$ is the orthogonal projection onto the $\mu$-eigenspace
of $\Delta_N$. Then on $\partial\Om\times[0,1]$,
\begin{equation}\label{eq:lateral}
    h(x,t)\big|_{\partial\Om}
    =\sum_{\mu\in\Sigma} e^{t(\lambda-\mu)/8}\,f_\mu(x).
\end{equation}

Assume the lateral parabolic Schiffer hypothesis
$h|_{\partial\Om\times[0,1]}\equiv c$. Then for every $x\in\partial\Om$,
\begin{equation}\label{eq:scalar}
    \sum_{\mu\in\Sigma} e^{t(\lambda-\mu)/8}\,f_\mu(x)=c
    \qquad\text{for all }t\in[0,1].
\end{equation}

\paragraph{Step 3: linear independence of the exponential profiles.}
The family $\{t\mapsto e^{t(\lambda-\mu)/8}:\mu\in\Sigma\}$ is
linearly independent in $C([0,1])$: distinct exponentials are
linearly independent (e.g.\ by repeated differentiation and the
non-vanishing Vandermonde). For each $x\in\partial\Om$,
\eqref{eq:scalar} thus forces
\begin{equation}\label{eq:percoeff}
    f_\mu(x)=\begin{cases} c & \mu=\lambda,\\
        0 & \mu\ne\lambda.\end{cases}
\end{equation}

Two sub-cases.

\emph{Case A: $\lambda\notin\Sigma$.} Then \eqref{eq:percoeff} reads
$f_\mu\equiv 0$ on $\partial\Om$ for every $\mu\in\Sigma$ and $c=0$.
Hence $\phi|_{\partial\Om}=\sum_\mu f_\mu\equiv 0$, and combined with
$\nor\phi=0$ on $\partial\Om$ this gives zero Cauchy data for the
second-order uniformly elliptic equation $L_{\Om,V}\phi=\lambda\phi$.
By Holmgren / Carleman unique continuation across a smooth boundary,
$\phi\equiv 0$ in $\Om$, contradicting non-triviality.

\emph{Case B: $\lambda\in\Sigma$.} Then $\lambda$ is a Neumann
eigenvalue of $-\Delta$ on $\Om$, $f_\lambda\equiv c$ on
$\partial\Om$, and $f_\mu\equiv 0$ on $\partial\Om$ for $\mu\ne\lambda$.

\paragraph{Step 4: extracting the classical Schiffer eigenfunction.}
In Case B, set $u:=P_\lambda\phi$. By construction, $u$ is a finite
linear combination of Neumann eigenfunctions of $-\Delta$ at
eigenvalue $\lambda$, hence
\[
    -\Delta u=\lambda u\quad\text{in }\Om,\qquad
    \nor u=0\quad\text{on }\partial\Om.
\]
Its boundary trace is $u|_{\partial\Om}=f_\lambda\equiv c$. Since
$c\ne 0$ (Case A is excluded), $u\not\equiv 0$. Thus $u$ is a
nontrivial classical Schiffer overdetermined eigenfunction on $\Om$,
proving the theorem.
\hfill$\square$

\section{Corollaries}

\begin{cor}[Radial test on the disc]\label{cor:disc}
Let $\Om=B\subset\R^n$ be the ball, $V\in C^\infty([0,R])$ a radial
convex potential. Then no smooth nontrivial pair
$(B,V,\phi,\lambda)$ with $V$ nonconstant satisfies lateral parabolic
Schiffer.
\end{cor}

\begin{proof}
For radial $V$, the Langevin operator preserves the radial sector and
the radial $\phi$ exists with $-\phi''-(n-1)r^{-1}\phi'+V'(r)\phi'
=\lambda\phi$ on $(0,R)$, $\phi'(R)=0$. The Schiffer overdetermination
$\phi(R)=c$ is automatic by radiality. By Theorem~\ref{thm:main}, if
lateral parabolic Schiffer holds, then $u=P_\lambda\phi$ is a
classical Schiffer eigenfunction of $-\Delta$ at $\lambda$ on $B$.
The classical Schiffer overdetermined eigenfunctions on the ball are
exactly the \emph{radial} Neumann eigenfunctions, with eigenvalue
$\lambda$ a radial Neumann eigenvalue of $-\Delta$ on $B$. So
$\lambda$ is a radial Neumann eigenvalue of $-\Delta$, AND $\lambda$
is the first Langevin eigenvalue of $L_{B,V}$. Subtracting the two
radial ODEs satisfied by $\phi$ (as a Langevin eigenfunction) and by
$u$ (as a $-\Delta$ eigenfunction): with both having the same
$\phi=u+\phi^\perp$ structure, the off-spectrum component $\phi^\perp$
satisfies zero Cauchy data on $\partial B$, hence vanishes by unique
continuation (it solves
$(-\Delta+V'\partial_r-\lambda)\phi^\perp=-V'u'$ with zero Cauchy
data). Thus $\phi=u$. But then
$L_{B,V}u=-\Delta u+V'(r)u'=\lambda u+V'(r)u'$, which equals
$\lambda u$ only if $V'(r)u'(r)\equiv 0$. Since $u$ is a nontrivial
real-analytic radial eigenfunction with $u'(R)=0$ but
$u'\not\equiv 0$, the zero set of $u'$ is discrete; hence $V'\equiv 0$,
i.e., $V$ is constant. Contradiction.
\end{proof}

\begin{cor}[Generic non-coincidence]\label{cor:generic}
If the Langevin first eigenvalue $\lambda_{\Om,V}$ is not a Neumann
eigenvalue of the unweighted Laplacian $-\Delta$ on $\Om$, then no
smooth nontrivial $\phi$ produces a lateral parabolic Schiffer triple
$(\Om,V,\phi,\lambda_{\Om,V})$. In particular, for any fixed convex
$\Om$ and any one-parameter family $\{V_s\}$ with $V_0\equiv 0$ and
$\partial_s V|_{s=0}\not\equiv 0$, the set of $s$ for which lateral
parabolic Schiffer can hold is discrete in $s$.
\end{cor}

\begin{proof}
The first statement is Case A of the theorem. The second follows
because $s\mapsto\lambda_{\Om,V_s}$ is real-analytic with nonzero
first-order shape derivative under generic perturbations, while the
Neumann spectrum of $-\Delta$ on $\Om$ is fixed and discrete; so
$\lambda_{\Om,V_s}$ matches an unweighted Neumann eigenvalue only on
a discrete set of $s$.
\end{proof}

\section{Reading the theorem from three angles}

\paragraph{Reduction.} Lateral parabolic Schiffer for weighted pairs
on $\Om$ reduces to (and is in fact equivalent up to the choice of
$V$ matching the eigenvalue, given Case B) classical Schiffer for the
unweighted Laplacian on $\Om$ at some eigenvalue. The potential $V$
adds no new geometric content; it can only resonate with an existing
Neumann eigenvalue of $-\Delta$.

\paragraph{Barrel dimension-reduction.} Combined with the cap
obstruction of the previous note (no barrel-lifted high-dim convex
Schiffer counterexample), the new theorem gives a sharper statement.
If a sequence of high-dimensional convex \emph{approximate} Schiffer
counterexamples $\Om_d$ has the barrel form $\Om_d=F_d(\Om,V)$, then
in the limit $d\to\infty$ they descend to a classical Schiffer
overdetermined eigenfunction on $\Om$. In particular, \emph{any
hypothetical high-dim convex Schiffer counterexample of barrel form
produces a classical Schiffer counterexample on the base $\Om$.}
There is no high-codimensional escape via the barrel route.

\paragraph{What about Direction 2 itself?} Direction 2 was: limits of
convex Schiffer counterexamples. Theorem~\ref{thm:main} narrows the
admissible limits among the enlarged class of weighted pairs
$(\Om,V)$. Specifically, the only way a sequence of weighted Schiffer
counterexamples can have a heat-extension parabolic-Schiffer
behaviour is if the underlying base $\Om$ already supports classical
Schiffer overdetermined data. So if classical Schiffer is known on a
class of base domains (low eigenvalues, centrally symmetric, etc.),
the corresponding weighted/parabolic Schiffer classes are
automatically empty over those bases.

\section{What this leaves open}

\begin{enumerate}[leftmargin=*]
    \item \textbf{The cap was the wrong target.} The reduction shows
    that the lateral parabolic Schiffer condition is automatically
    very rigid; the cap obstruction (Prop.~6 of the previous note) was
    the easy part. The real geometric content of weighted Schiffer is
    \emph{stationary} weighted Schiffer, which sits outside the
    lateral parabolic class. Concretely: a stationary weighted
    Schiffer pair $(\Om,V)$ at a Langevin eigenvalue $\lambda_{\Om,V}$
    not in the Neumann spectrum of $-\Delta$ exists without
    contradicting Theorem~\ref{thm:main}, but the heat extension of
    its eigenfunction simply does not behave well on the lateral
    boundary.

    \item \textbf{Stationary weighted Schiffer is still the right
    target for Strategy A.} The augmented compactness program of
    Strategy A of the original note now reads: enumerate stationary
    weighted Schiffer pairs $(\Om,V)$ and use \emph{stationary}
    overdetermined rigidity, not parabolic Serrin, to rule them out.
    Parabolic Serrin only applies to the (much smaller) subclass where
    $\lambda_{\Om,V}$ coincides with an unweighted Neumann eigenvalue.

    \item \textbf{The barrel construction has the wrong target.} The
    barrel limit produces an unweighted parabolic problem
    \eqref{eq:hseries}. To capture stationary weighted Schiffer the
    right limiting object would be the \emph{weighted} elliptic
    problem on $\Om$, which is what one would obtain via a warped
    barrel as suggested at the end of the previous note. Implementing
    this should be the next concrete construction.

    \item \textbf{A clean open question.} Does there exist a smooth
    convex pair $(\Om,V)$ with $V$ non-constant for which a stationary
    weighted Schiffer eigenfunction $\phi$ exists? Equivalently: does
    the weighted Pompeiu property hold for all $(\Om,V)$ with
    nonconstant $V$? This is the natural ``log-concave Schiffer
    conjecture'' aligned with the de Dios Pont log-concave extension
    of the hot spots conjecture.
\end{enumerate}

\end{document}
